\subsection{A certified simple central zero at exponent nineteen}
\label{subsec:nineteen-certificate}

We now certify the analytic hypothesis in the rank-one criterion.
Let $A$ be the genus-three quotient Jacobian, let $J$ be the
genus-nine Jacobi quotient Jacobian, and let $\chi$ be the order-nine
Dirichlet character modulo nineteen with $\chi(2)=\exp(2\pi i/9)$.
Write
\[
 L(A,s)=\sum_{n\geq1}a_n n^{-s},\qquad
 L_\chi(s)=L(A\otimes\chi,s),\qquad
 B=\left(\frac{19}{2\pi}\right)^3.
\]
The preceding Euler-factor calculation gives
$L(J,s)=L(A,s)L_\chi(s)L_{\bar\chi}(s)$.

\begin{proposition}[The central derivative is nonzero]
\label{prop:nineteen-analytic-certificate}
For the motivic normalization with center $s=1$,
\[
 2.4378<L'(A,1)<2.4383,\qquad
 4.087<|L_\chi(1)|^2<4.093,
\]
and
\[
 9.96<L'(J/\mathbb Q,1)<9.99.
\]
In particular $L(J/\mathbb Q,s)$ has a simple zero at one.
\end{proposition}

\begin{proof}
We give both truncation estimates, the finite exact input, and the
interval evaluation. Set $N=16384$. For $x>0$ define
\[
 H(x)=\frac1{2\pi i}\int_{\Re z=2}
          \frac{\Gamma(1+z)^3}{z}x^{-z}\,dz,
 \qquad
 D(x)=\frac1{2\pi i}\int_{\Re z=2}
          \frac{\Gamma(1+z)^3}{z^2}x^{-z}\,dz.
\]
If $X_1,X_2,X_3$ are independent exponential random variables of
mean one, Mellin inversion gives
\[
 H(x)=\Pr(X_1X_2X_3\geq x),\qquad
 D(x)=\mathbb E\!\left[\log_+\frac{X_1X_2X_3}{x}\right].
\]
Indeed $\mathbb E[(X_1X_2X_3)^z]=\Gamma(1+z)^3$;
the two inverse Mellin kernels are respectively the indicator of
$t>1$ and $\log_+t$. Thus, for every positive integer $k$,
\begin{equation}\label{eq:nineteen-kernel-bounds}
 0\leq H(x)\leq\Gamma(k+1)^3x^{-k},\qquad
 0\leq D(x)\leq\frac{\Gamma(k+1)^3}{k}x^{-k}.
\end{equation}
The first inequality is Markov's inequality. For the second, use
$\log_+t\leq t^k/k$ and take expectations.

Put $w=w(A\otimes\chi)$, the exact Gauss sum already calculated,
and
\[
 F_\chi=\sum_{n\geq1}\frac{a_n\chi(n)}n H(n/B).
\]
Contour shifting in the completed functional equations gives
\begin{equation}\label{eq:nineteen-central-sums}
 L_\chi(1)=F_\chi+w\overline{F_\chi},\qquad
 L'(A,1)=2\sum_{n\geq1}\frac{a_n}{n}D(n/B).
\end{equation}
For the first identity shift the contour of
$\Lambda_\chi(1+z)/z$, where
$\Lambda_\chi(s)=B^s\Gamma(s)^3L_\chi(s)$, and use
$\Lambda_\chi(s)=w\Lambda_{\bar\chi}(2-s)$.
The residue at zero is $\Lambda_\chi(1)$.
For the second use $\Lambda_A(1+z)/z^2$ and the sign
$w(A)=-1$; the residue is $\Lambda_A'(1)$.
Since $L(A,1)=0$ and $\Gamma(1)=1$, division by $B$
gives the stated derivative identity. The Hecke completed functions
are entire; their functional equations and gamma decay justify
the shifts. The kernels also give absolute convergence of the
resulting sums.

Let $b_n$ be the number of integral ideals of norm $n$ in the
degree-six CM field $K_0$. Its Hecke character has absolute value
$\sqrt{N\mathfrak a}$ on every ideal prime to nineteen. Grouping
the ideal series by norm gives
\[
 |a_n|\leq\sqrt n\,b_n,
 \qquad \sum b_n n^{-\sigma}=\zeta_{K_0}(\sigma)
 \quad(\sigma>1).
\]
Using \eqref{eq:nineteen-kernel-bounds}, the discarded tail of
$F_\chi$ is at most
\begin{equation}\label{eq:nineteen-coefficient-tail}
 E_N=\Gamma(k+1)^3B^k
       N^{\sigma-k-1/2}\zeta_{K_0}(\sigma).
\end{equation}
Here $k+1/2>\sigma$, and the inequality follows by extracting
$N^{\sigma-k-1/2}$ from each term with $n>N$.
The discarded derivative tail is at most $2E_N/k$.
We take $\sigma=11/10$ and $k=4$.

The zeta value in this bound is itself evaluated by an exact
Dirichlet factorization. If $\eta(2)=\exp(2\pi i/18)$ modulo
nineteen, then
\[
 \zeta_{K_0}(\sigma)=\zeta(\sigma)
     \prod_{j=1}^5 L(\eta^{3j},\sigma),
\]
and each nontrivial factor is the finite Hurwitz-zeta sum
\[
 L(\eta^k,\sigma)=19^{-\sigma}
   \sum_{j=0}^{17}\exp(2\pi i k j/18)
        \zeta\!\left(\sigma,\frac{2^j\bmod19}{19}\right).
\]
The factor $\zeta(\sigma)$ includes the prime above nineteen;
no Euler factor at that prime is omitted from the majorant.
Ball arithmetic gives
\[
 \zeta_{K_0}(11/10)<6.621273,
 \qquad E_N<0.000250861,
 \qquad 2E_N/4<0.000125431.
\]

For completeness, we also bound the infinite vertical contour
after the coefficient sum has been truncated. Put
$S_N=\sum_{n\leq N}|a_n|/n^3$ and $T=24$.
The recurrence for the gamma function and its Euler integral give
\[
 |\Gamma(3+it)|
   \leq\frac{\Gamma(3+m)}{|t|^m}\qquad(t\ne0).
\]
Consequently the discarded part of the half-sum contour for
$F_\chi$ is at most
\[
 V_H=\frac{B^2S_N\Gamma(3+m)^3}
             {\pi(3m)T^{3m}},
\]
while the discarded part of the derivative contour is at most
\[
 V_D=\frac{2B^2S_N\Gamma(3+m)^3}
             {\pi(3m+1)T^{3m+1}}.
\]
These estimates integrate $t^{-3m-1}$ and $t^{-3m-2}$,
respectively, on the two tails. With $m=16$ they give
$V_H<7.726\cdot10^{-19}$ and
$V_D<6.307\cdot10^{-20}$.

The finite input consists of all Euler factors needed through
$N$. Exactly 312 primes split completely in $K_0$ below this
cutoff. For each, an integral-basis vector of norm equal to the
prime is checked exactly; its three-conjugate product is checked
to have the correct Weil norm and residue sign, and its
characteristic polynomial supplies the Euler factor. This replay
uses only the number field and its maximal order, without a
class-group search or a class-number hypothesis. The remaining
residue-order cases use the explicit factors proved above and
the quadratic norm equation in $\mathbb Q(\sqrt{-19})$.
Inversion of these polynomials and multiplicativity give 468
nonzero coefficients through $N$. Each is an integer, and every
nonzero index is congruent to $1$, $7$ or $11$ modulo nineteen.

Here is the finite interval computation. For $z=2+it$ define
\[
 G(t)=B^z\Gamma(1+z)^3
       \sum_{n\leq N}\frac{a_n}{n^{1+z}}.
\]
The finite derivative is
$2\pi^{-1}\Re\int_0^T G(t)z^{-2}\,dt$.
For the complex half-sum, replace $a_n$ by $a_n\chi(n)$,
add its value at $t$ and at $-t$, and integrate with the
factor $(2\pi)^{-1}z^{-1}$ on the corresponding two terms.
This symmetric expression is meromorphic in the integration
variable and holomorphic near the real integration segments.
It does not apply complex conjugation to that variable
inside the interval integrator. The 24 unit intervals are
integrated at 128-bit precision by Arb, with requested absolute
and relative error tolerances $10^{-13}$. The returned enclosing
balls, rather than the requested tolerance alone, are used.
They give
\[
 \begin{split}
 D_N&\in[2.438030315780\mathbin\pm3.23\cdot10^{-13}],\\
 F_{\chi,N}&\in[0.9278931586\mathbin\pm3.40\cdot10^{-11}]\\
   &\hspace{12mm}+i[0.8339923732\mathbin\pm4.83\cdot10^{-11}].
 \end{split}
\]
The exact root number has enclosure
\[
 w\in[-0.9107472158164717387\mathbin\pm3\cdot10^{-20}]
       +i[0.4129642949245673588\mathbin\pm3\cdot10^{-20}].
\]
Adding the bounds $V_H+E_N$ to both coordinates of the
half-sum, adding $V_D+2E_N/4$ to the derivative, and then
using \eqref{eq:nineteen-central-sums} gives
\[
 \begin{split}
 L'(A,1)&\in[2.438\mathbin\pm0.000156],\\
 L_\chi(1)&\in[0.427\mathbin\pm0.000810]
                   +i[1.977\mathbin\pm0.000848],\\
 |L_\chi(1)|^2&\in[4.09\mathbin\pm0.00281],\\
 L'(A,1)|L_\chi(1)|^2&\in[9.97\mathbin\pm0.00891].
 \end{split}
\]
Finally, $L_{\bar\chi}(1)=\overline{L_\chi(1)}$ and
$L(A,1)=0$, so differentiating the factorization gives
$L'(J,1)=L'(A,1)|L_\chi(1)|^2$. The displayed enclosures
prove the inequalities and the simple zero.
\end{proof}

\begin{corollary}[The exponent-nineteen descent target]
\label{cor:nineteen-finite-sha}
The sixfold $P$ in Proposition~\ref{prop:lm-nineteen-rankone}
has rational rank zero and finite Tate--Shafarevich group. The groups
$\Sha(J/\Q)$ and $\Sha(A/\Q)$ are finite as well. Its generating
morphism $X_{19}\to P$ gives
$X_n(\A_\Q)_\bullet^{\mathrm{f\text{-}ab}}=X_n(\Q)$ for every
positive multiple $n$ of $19$.
\end{corollary}
\begin{proof}
Proposition~\ref{prop:nineteen-analytic-certificate} verifies the
analytic hypothesis of Proposition~\ref{prop:lm-nineteen-rankone}.
The latter gives full Tate--Shafarevich finiteness and a rational
rank-zero sixfold. Its proof constructs a map from $X_{19}$ whose
image generates $P$ by differences. The coordinate-power map
$X_n\to X_{19}$ therefore gives a nonconstant map to the same
rank-zero target with finite Tate--Shafarevich group.
\citet[Theorem~8.6, p.~385]{Stoll2007}, with its real-place erratum,
applies. This also proves Theorem~\ref{thm:nineteen-selmer-intro}.
\end{proof}
